\chapter{磁场}

电流周围为什么会出现环形的作用？为什么通电导线之间会互相吸引或排斥？为什么高速带电粒子进入磁场后会偏转成圆弧？这些问题共同指向电磁学中的第二种基本场：\emph{磁场}。与电场不同，磁场对静止电荷通常不施加作用，而是专门作用于\emph{运动电荷}与\emph{电流}。因此，磁现象天然地带有速度、方向与旋转的几何特征。

从微积分角度看，磁场一方面来源于连续分布的电流元，必须通过积分把每一个微元的贡献叠加起来；另一方面，带电粒子在磁场中的轨迹又满足微分方程
\[
  m\dv{\vb{v}}{t}=q\vb{v}\times\vb{B}.
\]
本章将围绕这两个核心问题展开：先用毕奥--萨伐尔定律计算磁场，再用洛伦兹力和安培力解释磁场中的运动与相互作用。

\section{磁场与磁感应强度}

在电场中，我们用电场强度 $\vb{E}$ 描述空间对电荷的作用；在磁场中，我们用磁感应强度 $\vb{B}$ 描述空间对运动电荷的磁作用。磁场的方向可以用小磁针北极受力方向定义，也可以用右手定则和电流方向联系起来。

\begin{definition}
磁场是由运动电荷、电流以及随时间变化的电场所激发的物理场。若把一个电荷量为 $q$、速度为 $\vb{v}$ 的试探粒子放入某处，磁场对它的作用由矢量 $\vb{B}$ 表征，称为\textbf{磁感应强度}。
\end{definition}

磁感应强度的国际单位制单位是特斯拉（T）。从量纲上看，$B$ 可以由力除以电荷速度得到，因此
\[
  1\,\mathrm{T}=1\,\mathrm{N\,s/(C\,m)}.
\]
在高中物理里，我们常把磁场画成磁感线；在大学电磁学里，磁感线只是帮助理解方向与疏密的几何图像，真正有物理意义的是每一点的矢量场 $\vb{B}(\vb{r})$。

\begin{law}[洛伦兹力]
电荷 $q$ 以速度 $\vb{v}$ 在磁场 $\vb{B}$ 中运动时，所受磁力为
\[
  \vb{F}_\mathrm{m}=q\vb{v}\times\vb{B}.
\]
其大小为
\[
  F_\mathrm{m}=|q|vB\sin\theta,
\]
其中 $\theta$ 是 $\vb{v}$ 与 $\vb{B}$ 的夹角。方向由叉乘决定：对正电荷用右手定则，对负电荷方向相反。
\end{law}

由上式立刻看出三种典型情形：
\begin{itemize}
  \item 当 $\theta=0$ 或 $\pi$ 时，$\vb{v}\parallel\vb{B}$，磁力为零；
  \item 当 $\theta=\pi/2$ 时，磁力最大，$F=|q|vB$；
  \item 磁力总垂直于速度，因此主要改变速度方向而不直接改变速率。
\end{itemize}

\begin{theorem}
纯磁场对带电粒子不做功。
\end{theorem}

\begin{proof}
磁力的瞬时功率为
\[
  P=\vb{F}_\mathrm{m}\cdot\vb{v}=q(\vb{v}\times\vb{B})\cdot\vb{v}=0,
\]
因为叉乘所得矢量恒与 $\vb{v}$ 垂直。因此磁场不能改变粒子的动能，只能改变其运动方向。
\end{proof}

这一定理极其重要。许多题目中粒子在磁场中做圆周运动、螺旋运动，速度大小保持不变，其根本原因都不是“圆周运动本来就匀速”，而是因为磁力不做功。

\begin{figure}[H]
\centering
\begin{tikzpicture}[scale=1.0, >=Stealth]
  \draw[thick] (0,-2.2) -- (0,2.2);
  \filldraw[black] (0,0) circle (0.06);
  \node[left] at (0,1.9) {$I$};
  \foreach \r in {0.6,1.1,1.6} {
    \draw[blue!70!black, thick] (0,0) circle (\r);
    \draw[blue!70!black, ->, thick] (\r,0) arc (0:55:\r);
  }
  \node[blue!70!black] at (2.1,1.5) {磁感线};
  \draw[->, red!70!black, thick] (2.8,0) -- (3.8,0) node[right] {$\vb{v}$};
  \filldraw[orange!80!black] (2.8,0) circle (0.07);
  \draw[->, violet!70!black, thick] (2.8,0) -- (2.8,1.1) node[above] {$\vb{F}$};
  \node at (2.8,-0.35) {$q>0$};
\end{tikzpicture}
\caption{长直电流周围的环形磁场与运动电荷受力示意}
\end{figure}

\begin{example}
一质量为 $m$、电荷量为 $q>0$ 的粒子以速度 $\vb{v}$ 射入匀强磁场，且 $\vb{v}$ 与 $\vb{B}$ 的夹角为 $30^\circ$。求磁力大小，并说明速度的哪一部分会被改变。
\end{example}

\begin{proof}
由洛伦兹力公式，
\[
  F= qvB\sin 30^\circ=\frac{1}{2}qvB.
\]
把速度分解为平行于磁场和垂直于磁场的两部分：
\[
  v_\parallel=v\cos30^\circ=\frac{\sqrt{3}}{2}v,
  \qquad
  v_\perp=v\sin30^\circ=\frac{1}{2}v.
\]
磁力始终垂直于 $\vb{B}$，也垂直于 $\vb{v}_\perp$，因此它只改变垂直分量的方向，而不改变平行分量。于是粒子最终将沿磁场方向前进，同时在垂直平面内做圆周运动，合成轨迹为螺旋线。
\end{proof}

\begin{remark}
许多初学者把“磁场方向”误认为“磁力方向”。实际上磁力方向同时取决于 $q$、$\vb{v}$ 和 $\vb{B}$ 三者；同一点的磁场方向固定，但不同速度方向的粒子在该点受力方向不同。
\end{remark}

\section{毕奥--萨伐尔定律}

既然磁场由电流激发，自然的问题是：给定一段导线和其中的电流，如何定量计算空间某点的 $\vb{B}$？对稳恒电流，这一问题的基本规律就是毕奥--萨伐尔定律。它与静电学中的库仑定律相似，都是从“微元贡献”出发，再对整个源分布积分。

设导线上一点处的长度微元为 $\dd\vb{\ell}$，电流为 $I$，场点相对该电流元的位置矢量为 $\vb{r}$，其长度为 $r$，对应单位矢量为 $\uvec{r}=\vb{r}/r$。则该电流元在场点产生的磁场微元满足：

\begin{law}[毕奥--萨伐尔定律]
稳恒电流元 $I\dd\vb{\ell}$ 在场点产生的磁感应强度微元为
\[
  \dd\vb{B}=\frac{\mu_0}{4\pi}\frac{I\,\dd\vb{\ell}\times\uvec{r}}{r^2}
  =\frac{\mu_0}{4\pi}\frac{I\,\dd\vb{\ell}\times\vb{r}}{r^3}.
\]
这里 $\mu_0$ 为真空磁导率，方向由 $\dd\vb{\ell}\times\uvec{r}$ 决定。
\end{law}

由叉乘大小公式，磁场微元的大小为
\[
  \dd B=\frac{\mu_0}{4\pi}\frac{I\,\dd\ell\sin\alpha}{r^2},
\]
其中 $\alpha$ 是电流元方向与指向场点方向之间的夹角。这说明：
\begin{itemize}
  \item 距离越远，贡献越小，呈 $1/r^2$ 衰减；
  \item 电流元与场点连线越接近平行，贡献越弱；
  \item 总磁场必须通过矢量叠加积分得到。
\end{itemize}

因此，对任意形状的细导线 $C$，磁场为
\[
  \vb{B}(\vb{r})=\frac{\mu_0 I}{4\pi}\int_C \frac{\dd\vb{\ell}'\times(\vb{r}-\vb{r}')}{|\vb{r}-\vb{r}'|^3}.
\]
这里 $\vb{r}'$ 是源点位置矢量，$\vb{r}$ 是场点位置矢量。这个积分表达式十分重要，它告诉我们：\emph{磁场是一个由电流分布生成的积分场}。

\begin{remark}
毕奥--萨伐尔定律本身已经隐含了叠加原理。若系统由多段导线组成，只需把每段导线的积分结果相加即可。这一点在求“圆弧加直线”“多匝线圈”“螺线管”等磁场时极为常用。
\end{remark}

为了熟悉微积分操作，我们先看一个有限直导线的统一积分模型。设导线沿 $x$ 轴放置，场点 $P$ 在 $y$ 轴上，距导线距离为 $a$。对位于坐标 $x$ 处的电流元，
\[
  \dd\vb{\ell}=\dd x\,\vu{x},
  \qquad
  \vb{r}=(-x)\vu{x}+a\vu{y},
  \qquad
  r=\sqrt{x^2+a^2}.
\]
于是
\[
  \dd B=\frac{\mu_0 I}{4\pi}\frac{a\,\dd x}{(x^2+a^2)^{3/2}},
\]
方向垂直纸面。积分就变成
\[
  B=\frac{\mu_0 I}{4\pi}\int_{x_1}^{x_2}\frac{a\,\dd x}{(x^2+a^2)^{3/2}}.
\]
注意到
\[
  \int \frac{a\,\dd x}{(x^2+a^2)^{3/2}}=\frac{x}{a\sqrt{x^2+a^2}}+C,
\]
因此有限直导线在点 $P$ 产生的磁场可写为
\[
  B=\frac{\mu_0 I}{4\pi a}\left(\frac{x_2}{\sqrt{x_2^2+a^2}}-\frac{x_1}{\sqrt{x_1^2+a^2}}\right).
\]
若引入两端点张角 $\theta_1,\theta_2$，则上式常写为
\[
  B=\frac{\mu_0 I}{4\pi a}(\sin\theta_1+\sin\theta_2).
\]
这说明几何角度与积分变量之间存在天然的三角替换关系。

\section{毕奥--萨伐尔定律应用}

本节用三个经典模型展示“电流元积分”的威力：长直导线、圆环以及长螺线管。它们既是电磁学中的基本结果，也是经典物理题中最常见的磁场背景。

\subsection*{一、无限长直导线的磁场}

无限长直导线可看成上节有限导线结果的极限，也可以直接积分。设导线沿 $x$ 轴无限延伸，场点 $P$ 到导线距离为 $a$。由上节得到
\[
  B=\frac{\mu_0 I}{4\pi}\int_{-\infty}^{\infty}\frac{a\,\dd x}{(x^2+a^2)^{3/2}}.
\]
令 $x=a\tan\phi$，则 $\dd x=a\sec^2\phi\,\dd\phi$，且
\[
  (x^2+a^2)^{3/2}=a^3\sec^3\phi.
\]
代入得
\[
  B=\frac{\mu_0 I}{4\pi a}\int_{-\pi/2}^{\pi/2}\cos\phi\,\dd\phi
  =\frac{\mu_0 I}{4\pi a}\bigl[\sin\phi\bigr]_{-\pi/2}^{\pi/2}
  =\frac{\mu_0 I}{2\pi a}.
\]
方向由右手定则给出，沿绕导线的切向。

\begin{example}
证明无限长直导线周围磁场满足 $B\propto 1/a$，并说明为什么离导线越远磁感线越稀疏。
\end{example}

\begin{proof}
由积分结果
\[
  B(a)=\frac{\mu_0 I}{2\pi a}
\]
知，在电流 $I$ 固定时，磁场大小与距离 $a$ 成反比。几何上，磁感线是以导线为中心的同心圆。若以半径 $a$ 的圆周作为“绕线一圈”的几何长度，则同样数量级的环量分布在更长的圆周上，单位长度分配到的磁场自然减小，所以图像上表现为越远处磁感线越稀疏。
\end{proof}

\begin{figure}[H]
\centering
\begin{tikzpicture}[scale=1.0, >=Stealth]
  \draw[thick] (0,-2.2) -- (0,2.2);
  \node[left] at (0,2.0) {$I$};
  \filldraw[black] (0,0) circle (0.07);
  \draw[red!70!black, thick] (3,0) circle (0.06);
  \node[right] at (3,0) {$P$};
  \draw[dashed] (0,0) -- (3,0) node[midway, above] {$a$};
  \foreach \r in {0.8,1.4,2.0} {
    \draw[blue!70!black, thick] (0,0) circle (\r);
    \draw[blue!70!black, ->, thick] (\r,0) arc (0:40:\r);
  }
  \node[blue!70!black] at (2.1,1.55) {$\vb{B}$ 沿圆周切向};
\end{tikzpicture}
\caption{无限长直导线的环形磁场}
\end{figure}

\subsection*{二、载流圆环轴线上的磁场}

设半径为 $R$ 的圆环位于 $xy$ 平面，电流为 $I$，求其轴线上距圆心 $z$ 处点 $P$ 的磁场。由对称性，圆周上相对的电流元在横向分量上互相抵消，只剩轴向分量相加。

对任意电流元，场点距离均为
\[
  r=\sqrt{R^2+z^2}.
\]
由于 $\dd\vb{\ell}\perp\vb{r}$，微元磁场大小为
\[
  \dd B=\frac{\mu_0 I}{4\pi}\frac{\dd\ell}{r^2}.
\]
其轴向分量为
\[
  \dd B_z=\dd B\sin\beta
  =\frac{\mu_0 I}{4\pi}\frac{\dd\ell}{r^2}\cdot\frac{R}{r}
  =\frac{\mu_0 I R}{4\pi r^3}\dd\ell,
\]
其中 $\sin\beta=R/r$。沿整个圆环积分：
\[
  B_z=\frac{\mu_0 I R}{4\pi r^3}\int \dd\ell
  =\frac{\mu_0 I R}{4\pi (R^2+z^2)^{3/2}}\cdot 2\pi R.
\]
故
\[
  B(z)=\frac{\mu_0 I R^2}{2(R^2+z^2)^{3/2}}.
\]

\begin{figure}[htbp]
\centering
\begin{tikzpicture}
\begin{axis}[
  width=0.75\textwidth, height=0.5\textwidth,
  xlabel={轴向位置 $z/R$}, ylabel={磁场 $B/B_0$},
  axis lines=middle,
  xmin=-3.2, xmax=3.2, ymin=0, ymax=1.08,
  grid=major, grid style={gray!30},
  thick,
]
\addplot[blue, domain=-3:3, samples=220] {(1 + x^2)^(-1.5)};
\end{axis}
\end{tikzpicture}
\caption{圆形电流环轴线上磁场在圆心处最强并向两侧衰减}
\end{figure}

特别地，在圆心 $z=0$ 处，
\[
  B_\mathrm{center}=\frac{\mu_0 I}{2R}.
\]

\begin{example}
半径为 $R$、电流为 $I$ 的单匝圆环，在轴线上距圆心 $z=\sqrt{3}R$ 处的磁感应强度是多少？
\end{example}

\begin{proof}
代入轴线公式：
\[
  B=\frac{\mu_0 I R^2}{2(R^2+3R^2)^{3/2}}
  =\frac{\mu_0 I R^2}{2(4R^2)^{3/2}}
  =\frac{\mu_0 I R^2}{2\cdot 8R^3}
  =\frac{\mu_0 I}{16R}.
\]
这个结果表明，轴线上磁场随距离衰减很快，远处呈近似偶极子场的趋势。
\end{proof}

\begin{figure}[H]
\centering
\begin{tikzpicture}[scale=1.0, >=Stealth]
  \draw[thick] (0,0) ellipse (2.3 and 0.8);
  \draw[dashed] (-2.3,0) -- (2.3,0);
  \draw[->, red!70!black, thick] (0,-1.8) -- (0,2.2) node[above] {$z$};
  \filldraw[blue!70!black] (0,1.5) circle (0.06);
  \node[right] at (0,1.5) {$P(z)$};
  \draw[dashed] (0,0) -- (0,1.5) node[midway, right] {$z$};
  \draw[dashed] (0,0) -- (2.3,0) node[midway, above] {$R$};
  \draw[->, violet!70!black, thick] (0,1.5) -- (0,2.0) node[right] {$\vb{B}$};
  \draw[->, thick] (1.5,0.53) -- (1.1,0.72);
\end{tikzpicture}
\caption{载流圆环轴线上的磁场方向}
\end{figure}

\subsection*{三、螺线管磁场的积分估计}

把长螺线管看成许多相邻圆环沿轴线连续堆叠，每单位长度匝数为 $n$，每匝电流为 $I$。设观察点位于螺线管中心附近。对位于轴坐标 $\zeta$ 的薄圆环，其在中心点产生的轴向磁场微元为
\[
  \dd B=\frac{\mu_0 I R^2}{2(R^2+\zeta^2)^{3/2}}\,n\dd\zeta.
\]
若螺线管足够长，可把积分区间近似视为 $(-\infty,\infty)$，于是
\[
  B\approx \frac{\mu_0 n I R^2}{2}\int_{-\infty}^{\infty}\frac{\dd\zeta}{(R^2+\zeta^2)^{3/2}}.
\]
利用积分公式
\[
  \int \frac{\dd\zeta}{(R^2+\zeta^2)^{3/2}}=\frac{\zeta}{R^2\sqrt{R^2+\zeta^2}}+C,
\]
得到
\[
  B\approx \frac{\mu_0 n I R^2}{2}\cdot \frac{2}{R^2}=\mu_0 n I.
\]
这就是长直螺线管内部磁场近似均匀的著名结果。

\begin{example}
一个长螺线管每单位长度有 $n$ 匝，通有电流 $I$。试用圆环叠加积分说明其中心附近磁场近似与半径无关，只与 $nI$ 成正比。
\end{example}

\begin{proof}
由上式可知，中心磁场表达式中电流和匝密度只以乘积 $nI$ 出现；而积分
\[
  \int_{-\infty}^{\infty}\frac{R^2\dd\zeta}{(R^2+\zeta^2)^{3/2}}=2
\]
是一个与 $R$ 无关的纯数。也就是说，半径增大虽然使单匝在中心的磁场减弱，但能贡献的几何范围同步增大，两种效应相互抵消，最终中心附近得到
\[
  B\approx \mu_0 nI.
\]
这正是“长螺线管内部磁场主要由轴向匝密度决定”的微积分解释。
\end{proof}

\begin{remark}
毕奥--萨伐尔定律总能用于计算磁场，但积分往往繁琐。对高对称问题，更高效的方法是下一节的安培环路定理。不过，安培环路定理只能在对称性足够强时快速求值；若缺乏对称性，毕奥--萨伐尔积分依然是基本武器。
\end{remark}

\section{安培环路定理}

对于圆柱对称、平面对称或环形对称的电流分布，毕奥--萨伐尔定律虽然可用，但往往不如环路积分简洁。安培环路定理把“电流产生磁场”这一关系写成沿闭合曲线的积分形式。

\begin{law}[安培环路定理]
在真空中的稳恒电流情形下，沿任意闭合曲线 $C$ 的磁场环量满足
\[
  \oint_C \vb{B}\cdot\dd\vb{\ell}=\mu_0 I_{\mathrm{enc}},
\]
其中 $I_{\mathrm{enc}}$ 表示闭合曲线所围曲面所穿过的代数电流。
\end{law}

这一定理告诉我们：闭合回路上磁场的切向分量积分，只取决于所包围的总电流，而与所选曲面的具体形状无关。在数学上，它与斯托克斯公式和微分形式 $\curl\vb{B}=\mu_0\vb{J}$ 密切相关；在高中阶段，我们重点掌握其对称性用法。

\subsection*{一、无限长直导线再求磁场}

对半径为 $r$、以导线为中心的圆形安培回路，由对称性可知回路上各点磁场大小相等、方向沿切线，因此
\[
  \oint \vb{B}\cdot\dd\vb{\ell}=B\oint \dd\ell=B(2\pi r).
\]
包围电流为 $I$，故
\[
  B(2\pi r)=\mu_0 I,
\qquad
  B=\frac{\mu_0 I}{2\pi r}.
\]
这个结果与毕奥--萨伐尔积分完全一致，但推导长度大大缩短。

\subsection*{二、长直螺线管内部磁场}

取一条长方形安培回路，其中一边位于螺线管内部并平行于轴线，长度为 $\ell$；另一边位于管外。对足够长的螺线管，外部磁场近似为零，两条短边又与磁场垂直，因此环量近似仅来自内部那一边：
\[
  \oint \vb{B}\cdot\dd\vb{\ell}\approx B\ell.
\]
若单位长度匝数为 $n$，则回路包围的总电流为 $n\ell\,I$，于是
\[
  B\ell=\mu_0 n\ell I,
  \qquad
  B=\mu_0 nI.
\]
这里我们看到，安培环路定理把原本复杂的叠加积分压缩成了一个几何量 $n\ell$。

\subsection*{三、环形螺线管（环形铁心线圈）}

设一环形线圈总匝数为 $N$，通电流 $I$，求距圆心 $r$ 处的磁场。取与环形线圈同心的圆形安培回路，则
\[
  B(2\pi r)=\mu_0 NI,
\]
故
\[
  B(r)=\frac{\mu_0 NI}{2\pi r}.
\]
可见环形线圈内部磁场沿环向分布，且随半径呈 $1/r$ 变化；在线圈外部，由于包围净电流为零，理想情况下磁场近似为零。

\begin{theorem}
安培环路定理在求磁场大小时，只有在能够事先判断回路上 $B$ 的大小恒定或某些线段贡献为零时，才具有直接的计算价值。
\end{theorem}

\begin{proof}
由
\[
  \oint_C \vb{B}\cdot\dd\vb{\ell}=\mu_0 I_{\mathrm{enc}}
\]
只能得到环量信息，而不能直接给出每一点的 $B$。若因对称性可断定：
\begin{enumerate}
  \item 在某段回路上 $\vb{B}$ 与 $\dd\vb{\ell}$ 平行且大小恒定；
  \item 在某段回路上 $\vb{B}\perp\dd\vb{\ell}$；
  \item 某些区域 $B\approx 0$；
\end{enumerate}
则环量积分才能简化为代数乘积，从而解出 $B$。若不具备这些条件，环路定理仍然成立，但不足以单独求解局部磁场分布。
\end{proof}

\begin{remark}
因此，毕奥--萨伐尔定律是“通用公式”，安培环路定理是“对称性捷径”。掌握何时该积分、何时该选环路，是磁场问题中的关键能力。
\end{remark}

\section{带电粒子在磁场中运动}

磁场对运动电荷施加始终垂直于速度的力，因此带电粒子的轨迹通常是曲线。只要磁场足够均匀，运动规律可以由牛顿第二定律与洛伦兹力方程直接写出。

\subsection*{一、速度垂直于磁场：匀速圆周运动}

若 $\vb{v}\perp\vb{B}$，则磁力大小恒为 $|q|vB$，始终指向圆心方向，恰好充当向心力：
\[
  |q|vB=\frac{mv^2}{r}.
\]
因此轨道半径为
\[
  r=\frac{mv}{|q|B}.
\]
角速度为
\[
  \omega=\frac{v}{r}=\frac{|q|B}{m},
\]
周期为
\[
  T=\frac{2\pi}{\omega}=\frac{2\pi m}{|q|B}.
\]
值得注意的是，在非相对论范围内，周期与速率无关，只与 $m/(|q|B)$ 有关，这为回旋加速器、质谱仪等装置提供了理论基础。

\begin{figure}[htbp]
\centering
\begin{tikzpicture}
\begin{axis}[
  width=0.75\textwidth, height=0.5\textwidth,
  xlabel={速度 $v$}, ylabel={轨道半径 $r$},
  axis lines=middle,
  xmin=0, xmax=5.2, ymin=0, ymax=3.2,
  grid=major, grid style={gray!30},
  thick,
]
\addplot[blue, domain=0:5, samples=120] {0.6*x};
\end{axis}
\end{tikzpicture}
\caption{在 $m$、$q$ 与 $B$ 固定时，圆周运动半径与粒子速度成正比}
\end{figure}

\subsection*{二、速度斜射于磁场：螺旋运动}

若初速度与磁场夹角为 $\theta$，则分解为
\[
  v_\parallel=v\cos\theta,
  \qquad
  v_\perp=v\sin\theta.
\]
其中 $v_\perp$ 导致圆周运动，半径为
\[
  r=\frac{mv_\perp}{|q|B};
\]
而 $v_\parallel$ 不受磁力影响，保持匀速直线运动。因此总轨迹是沿磁场方向前进的螺旋线。其螺距（相邻两圈轴向距离）为
\[
  h=v_\parallel T=\frac{2\pi m v_\parallel}{|q|B}.
\]

\begin{theorem}
在匀强磁场中，带电粒子的速度大小保持不变；若初速度分解为 $\vb{v}=\vb{v}_\parallel+\vb{v}_\perp$，则轨迹必为圆、直线或螺旋线中的一种。
\end{theorem}

\begin{proof}
由上一节结论，磁场不做功，因此速率 $v$ 恒定。若 $v_\perp=0$，粒子不受力，做直线运动；若 $v_\parallel=0$，粒子只做圆周运动；若二者都不为零，则垂直分量负责圆周运动，平行分量负责匀速前进，叠加后得到螺旋线。
\end{proof}

\begin{example}
一个电子以速度 $v$ 射入磁感应强度为 $B$ 的匀强磁场，速度与磁场夹角为 $60^\circ$。求其轨道半径、周期和螺距。
\end{example}

\begin{proof}
电子电荷量绝对值记为 $e$。有
\[
  v_\perp=v\sin60^\circ=\frac{\sqrt{3}}{2}v,
  \qquad
  v_\parallel=v\cos60^\circ=\frac{1}{2}v.
\]
故轨道半径
\[
  r=\frac{mv_\perp}{eB}=\frac{\sqrt{3}mv}{2eB}.
\]
周期
\[
  T=\frac{2\pi m}{eB}.
\]
螺距
\[
  h=v_\parallel T=\frac{1}{2}v\cdot\frac{2\pi m}{eB}=\frac{\pi mv}{eB}.
\]
由于电子带负电，其旋转方向与正电荷按右手定则得到的方向相反，但上述半径和周期的大小不受符号影响。
\end{proof}

\begin{figure}[H]
\centering
\begin{tikzpicture}[scale=0.95, >=Stealth]
  \draw[->, thick] (0,-0.2) -- (0,4.5) node[above] {$\vb{B}$};
  \draw[blue!70!black, thick, samples=120, domain=0:10, smooth, variable=\t]
    plot ({0.7*sin(360*\t)}, {0.4*\t});
  \draw[->, red!70!black, thick] (0.2,0.3) -- (1.3,0.9) node[right] {$\vb{v}$};
  \node[blue!70!black] at (1.7,3.6) {螺旋轨迹的二维投影};
\end{tikzpicture}
\caption{带电粒子斜射入匀强磁场时的螺旋运动示意}
\end{figure}

\begin{remark}
若题目给出的是有限磁场区域，关键常常不在“套公式”，而在几何关系：圆心位置、半径与边界的关系决定粒子从哪里射出。微积分负责给出半径和周期，几何负责把它转化为位移、时间和偏转角。
\end{remark}

\section{安培力}

磁场不仅作用于单个运动电荷，也作用于宏观电流。导线中大量带电粒子定向移动，磁场对它们的总作用可以等效为对整段导线的作用，这就是安培力。

\begin{law}[安培力公式]
长度矢量为 $\vb{L}$、电流为 $I$ 的直导线置于匀强磁场 $\vb{B}$ 中时，所受安培力为
\[
  \vb{F}=I\vb{L}\times\vb{B}.
\]
其大小为
\[
  F=ILB\sin\theta,
\]
其中 $\theta$ 为电流方向与磁场方向夹角。
\end{law}

这一公式可由洛伦兹力从微观上理解。设导线横截面积为 $S$，载流子数密度为 $n$，每个载流子电荷量为 $q$，定向漂移速率为 $v_\mathrm{d}$。长度为 $L$ 的一段导线内载流子总数为 $nSL$，总受力为
\[
  F=nSL\cdot qv_\mathrm{d}B.
\]
而电流满足
\[
  I=nqSv_\mathrm{d},
\]
故
\[
  F=ILB.
\]
若方向不垂直，则取叉乘形式即可。这说明安培力本质上是无数个载流子所受洛伦兹力的宏观总和。

\subsection*{平行导线间的相互作用}

设两条相距 $d$ 的长直平行导线，分别通电流 $I_1$、$I_2$。第一条导线在第二条处产生磁场
\[
  B_1=\frac{\mu_0 I_1}{2\pi d}.
\]
于是第二条导线单位长度所受安培力为
\[
  \frac{F}{L}=I_2B_1=\frac{\mu_0 I_1 I_2}{2\pi d}.
\]
若两电流同向，则两导线相互吸引；若反向，则相互排斥。这个结论是“电流之间存在磁相互作用”的最典型体现。

\begin{example}
两条无限长平行直导线相距 $d$，通有同向电流 $I_1$ 和 $I_2$。求每单位长度相互作用力，并说明其方向。
\end{example}

\begin{proof}
由第一条导线在第二条导线处产生的磁场大小
\[
  B_1=\frac{\mu_0 I_1}{2\pi d},
\]
第二条导线单位长度受力为
\[
  \frac{F_{21}}{L}=I_2B_1=\frac{\mu_0 I_1I_2}{2\pi d}.
\]
方向由 $\vb{L}\times\vb{B}$ 决定，指向第一条导线；同理，第一条导线也受大小相同、方向相反的力。故两导线相互吸引。若其中一条电流反向，则叉乘方向改变，相互作用变为排斥。
\end{proof}

\begin{remark}
安培力是工程中电机、电流天平、扬声器等装置的基础。其宏观公式简洁，但使用时应先判断磁场是否匀强；若磁场随位置变化，就需要把导线分成微元并积分：
\[
  \vb{F}=I\int \dd\vb{\ell}\times\vb{B}.
\]
这与毕奥--萨伐尔定律中“由电流元求磁场”的思想互为镜像。
\end{remark}

下面选取三类典型经典题型：粒子偏转、速度选择和导线受力。题目的物理背景属于高中范围，但我们用微积分视角把公式之间的联系讲得更清楚。

\section*{本章小结}

本章的逻辑可以概括为四步：
\begin{enumerate}
  \item 用磁感应强度 $\vb{B}$ 描述磁场，用洛伦兹力 $q\vb{v}\times\vb{B}$ 描述其对运动电荷的作用；
  \item 用毕奥--萨伐尔定律把电流元的磁场贡献积分起来，得到长直导线、圆环和螺线管等经典结果；
  \item 用安培环路定理在高对称情形下快速求场，理解环量与包围电流之间的关系；
  \item 用牛顿第二定律分析带电粒子轨迹，并把洛伦兹力推广为电流元受力，得到安培力与平行导线间相互作用。
\end{enumerate}

从微积分角度说，磁学问题常在“积分求场”和“微分方程求轨道”之间来回切换。看到电流分布，就想到把微元叠加；看到粒子轨迹，就想到把受力方程写成曲率或周期关系。这正是“用微积分重新理解高中物理”的核心训练。

